\question\TFQuestion{T}{If $p^2 = 0.16$, then total homozygosity in the population is 0.52.}
\question\TFQuestion{F}{Gene flow is enough to stop the effects of genetic drift in a large population.}
\question\TFQuestion{F}{One difference between intrasexual and intersexual selection is that intrasexual selection does not cause sexual dimorphism.}
\question\TFQuestion{T}{Genetic variation among individuals in a population is necessary for natural selection to act upon individuals in that population.}
\question\TFQuestion{T}{Microevolution and macroevolution are caused by the same underlying process: change in genetic variation in a population over time.}
\question\TFQuestion{F}{If a trait is energetically costly and provides no survival advantage, then genetic drift will cause that trait to disappear from the population over many generations.}
\question\TFQuestion{T}{Sexual dimorphism is most often the result of directional selection.}
\question\TFQuestion{F}{One assumption of the Hardy-Weinberg equations is that all individuals have an equally likely chance of migrating to another population.}
\question\TFQuestion{F}{Overproduction of offspring by every individual in a population is necessary for evolution to occur.}
\question\TFQuestion{F}{Natural selection by genetic drift is most common in small populations.}
\question\TFQuestion{F}{Heterochrony can be defined as the retention of juvenile characteristics in the adult.}
\question\TFQuestion{F}{Adaptations cause descent with modification.}
\question\TFQuestion{F}{The smallest unit of evolution is the individual.}
\question\TFQuestion{F}{If populations of a single species are geographically isolated for any period of time, then they will evolve biological reproductive isolation.}